% sec5-periodicity.tex -- Section 5 of the C3 paper (draft 2026-11-08): result E.
% Sources (run B 20261104-0759): claude-20261104-c3-closure/01-PC-CLOSURE-PERIODICITY.md (c1: closure formulation,
% tilting proof, R1-R3, computations); codex-20261104-c3-monoidal/05-CONSOLIDATED-EXACT-PC-PERIODICITY.md (x2: primary
% proof), 03 (equality-boundary counterexample), 04 (exact-quotient repair); codex-20261104-x1-adjudication/01;
% notes-20261104/referee-runB-closure-periodicity.md (cold referee); Astra codex-20261104-runB-adjudication/00 (C1) and 03.
% Hone et al. statements checked against the archived arXiv:2509.07377v2 text (Def 6.1.1, Lemma 7.3.2, Thm 8.4.5,
% Thm 8.5.3, Lemma 8.5.1) on 2026-11-08.

\section{Periodicity and the exact degree}\label{sec:periodicity}

This section proves result E: whether $\mathrm{onto\text{-}FPHP}^{m}_n$ has a PC refutation of
degree at most $D$ depends, for $m$ large compared with $D$, only on $m$ modulo the least
power of $p$ above $D$.  Since the square board $m=n$ is satisfiable, this fixes the PC degree
exactly.

\begin{theorem}[periodicity]\label{thm:periodicity}
Let $F$ be a field of characteristic $p$, $n\ge1$, $D\ge1$, and $q$ a power of $p$ with
$q>D$.  Let $m\equiv m'\pmod q$.
\begin{enumerate}
\item[(a)] If $m,m'>2D$, then $\mathrm{onto\text{-}FPHP}^{m}_n$ has a $\mathrm{PC}_F$
  refutation of degree at most $D$ if and only if $\mathrm{onto\text{-}FPHP}^{m'}_n$ has one.
\item[(b)] If Theorem~8.4.5 of \cite{hone2025} holds (see below), the same holds for
  $m,m'\ge2D$.
\end{enumerate}
\end{theorem}

Here $m$ is the number of pigeons and $n$ the number of holes; $n$ is arbitrary, and $m<n$ is
allowed.

\begin{corollary}[E, exact degree]\label{cor:exact}
Let $p$ be prime, $k\ge1$, $n\ge1$ and $q=p^{v_p(k)}$.  Over every field of characteristic
$p$, the PC degree of $\mathrm{onto\text{-}FPHP}^{n+k}_n$ is exactly $q$ for every
$n\ge2q-1$.  If Theorem~8.4.5 of \cite{hone2025} holds, the same holds for $n\ge2q-2$.
\end{corollary}

\begin{proof}
If $q=1$ the degree is $1$ (Section~\ref{sec:intro-question}).  Let $q\ge p$ and $D=q-1$.
The board with $n$ pigeons and $n$ holes has a solution, a permutation matrix, so it has no
refutation of any degree.  Since $n\equiv n+k\pmod q$ and $n>2D$ when $n\ge2q-1$,
Theorem~\ref{thm:periodicity}(a) gives no refutation of degree at most $q-1$ for $n+k$
pigeons.  Proposition~\ref{prop:riis} gives one of degree $q$.  With (b), $n\ge2D=2q-2$
suffices.
\end{proof}

In particular, for surplus $p^\ell$ the PC degree is exactly $p^\ell$ once $n\ge2p^\ell-1$:
Riis's upper bound is attained, and the degree grows with $\ell$ as fast as it can.  This is
the strongest statement in this paper.  It is also the one that rests on the deepest
inputs, so we state its status with care.

\paragraph{Tier and inputs.}  T-WRITTEN.  Part (a) is proved in Appendix~\ref{app:tilting}
(outlined in Section~\ref{sec:periodicity-tilting}), using Harman's
equivalences~\cite{harman2015} with the repairs of Section~\ref{sec:errata}, and standard facts
on tilting modules (tensor products and Levi restrictions of tilting modules are tilting;
bounded tilting complexes compute the derived category).  The input we consider most exposed
there is the compatibility of Harman's equivalence with induction and restriction of
arbitrary submodules (Lemma~\ref{lem:B}), whose tilting-theoretic foundations we cite and do
not re-prove.

Section~\ref{sec:periodicity-oti} gives a shorter argument through the semisimplification
functors of Hone, Klein, Pauwels, Sherman, Yacobi and Zhang~\cite{hone2025}, which covers (b)
as well.  It uses their Theorem~8.4.5.  A model reconstruction of the proof of that theorem
(Astra, November 2026) found that the published proof does not establish it: their
Lemma~8.4.4 rests on Lemma~8.2.1, which is false as stated (Section~\ref{sec:errata}).  This is
not a counterexample to Theorem~8.4.5, and the reconstruction found every other step of
Section~\ref{sec:periodicity-oti}, and their Lemma~7.3.2, to hold.  Until that step is
repaired, (b) and the threshold $n\ge2q-2$ are conditional, and (a) rests on
Appendix~\ref{app:tilting}.

\paragraph{Comparison.}  For $q=p$, Theorem~A gives the exact degree $p$ at the Lean tier
from $n\ge2p-3$.  For $q=p^2$, result F gives the exact degree $p^2$ from $n\ge2p^2-3$ for
every prime, on different inputs (Section~\ref{sec:bands}).  Result D gives $2^\ell$, which is
exact when $p=2$, on the inputs and thresholds of Section~\ref{sec:growth}.
Below the threshold the degree can be smaller: over $\mathbb F_2$ with $k=4$, exact
computation of the closure below gives degree $3$ at $n=3$ and $n=4$, while Lean gives
exactly $4$ from $n=5$ (\texttt{[1156] StdPCExact.std\_f2\_k4\_exact}), below the threshold
$2q-1=7$ of Corollary~\ref{cor:exact}.

\subsection{The closure}\label{sec:periodicity-closure}

We use the notation of Section~\ref{sec:defs-written} on the board with $m$ pigeons and $n$
holes, and write $\mathcal B_m$, $F_d(m)$.  Let $C_0(m)\subseteq F_D(m)$ be the span of the
reduced pigeon and hole axioms, and recursively
\[
  S_t=C_t\cap F_{D-1},\qquad
  C_{t+1}=C_t+\sum_{i,j}\rho(x_{ij})\,S_t .
\]
The chain is increasing in the finite-dimensional space $F_D(m)$, so it stabilises; let
$C_\infty(m)$ be its limit.  All $C_t$ are invariant under permutations of the pigeons.

\begin{lemma}\label{lem:closure}
$\mathrm{onto\text{-}FPHP}^m_n$ has a $\mathrm{PC}_F$ refutation of degree at most $D$ if and
only if $1\in C_\infty(m)$.
\end{lemma}

\begin{proof}
($\Rightarrow$) As in Proposition~\ref{prop:nd-pc}, every line $f$ of a degree-$D$
derivation has $\rho(f)\in C_\infty$: axioms reduce into $C_0$ or to $0$, linear combinations
stay in $C_\infty$, and if $xf$ is a legal multiplication then $\deg f\le D-1$, so
$\rho(f)\in C_\infty\cap F_{D-1}=S_t$ for large $t$ and $\rho(xf)\in C_{t+1}$.

($\Leftarrow$) We show by induction on $t$ that every element of $C_t$, written as a
multilinear matching polynomial, is a line of a degree-$D$ derivation.  For $t=0$ this is
clear.  Let $s\in S_t$; by induction $s$ is derivable, and $\deg s\le D-1$ as a polynomial,
so $x_{ij}s$ is a legal multiplication.  Its reduction $\rho(x_{ij}s)$ differs from
$x_{ij}s$ by multiples of Boolean, collision and functional axioms: reduce one repeated
variable or forbidden pair at a time; each multiple has degree at most $\deg(x_{ij}s)\le D$,
as in Lemma~\ref{lem:ml-reduce}.  Linear combinations give $C_{t+1}$.
\end{proof}

On the square board ($m=n$), every element of $C_\infty$ vanishes at every permutation
matrix (true for $C_0$ and preserved by the recursion), so $1\notin C_\infty(n)$.

The recursion is a finite sequence of linear-algebra operations on matrices with entries in
$\mathbb F_p$.  Extending scalars to any field containing $\mathbb F_p$ commutes with sums,
images and intersections, so $1\in C_\infty$ over $F$ if and only if $1\in C_\infty$ over
$\mathbb F_p$ if and only if over $\overline{\mathbb F}_p$.  We may therefore work over
$k=\overline{\mathbb F}_p$.

\paragraph{The closure as representation theory.}  The symmetric group $S_m$ acts on
pigeons, and all spaces above are $kS_m$-modules.  For $e\ge0$ and a fixed set of $e$ holes
listed in increasing order, the $e$-matchings onto those holes form the permutation module
$E_e(m)$ of ordered $e$-tuples of distinct pigeons; so $F_d(m)$ is a direct sum of
$\binom ne$ copies of $E_e(m)$, $e\le d$.  For a hole $j$ let $U_j=E_1(m)$, the natural
permutation module on pigeons, and let $\mu_j\colon U_j\otimes F_{D-1}(m)\to F_D(m)$ send
$u_i\otimes x_M$ to $\rho(x_{ij}x_M)$.  Then
\[
  C_{t+1}=C_t+\textstyle\sum_{j}\mu_j\bigl(U_j\otimes S_t\bigr),
\]
and the constant $1$ is the image of the map $k\to F_D(m)$ onto the empty matching.  So
$C_\infty(m)$ is built from permutation modules by tensoring with $U_j$, images,
intersections and sums, and the question is whether the composite
$k\to F_D(m)\to F_D(m)/C_\infty(m)$ is zero.

\subsection{A conditional proof through semisimplification}\label{sec:periodicity-oti}

This argument is shorter than the one in Appendix~\ref{app:tilting} and covers part (b), but
it depends on \cite[Thm.~8.4.5]{hone2025}, whose published proof is incomplete
(Section~\ref{sec:errata}).  We give it because every other step has been checked, so a
repair of that one step would make (b) unconditional.

It suffices to prove one step: $M=L+q$ pigeons against $L$ pigeons, with $L\ge2D$ (for (b))
or $L>2D$ (for (a)).  Applying the step repeatedly in either direction gives the theorem.

\paragraph{The functor.}  Let $H\subseteq S_q$ be an elementary abelian $p$-group of order
$q$ acting regularly on the last $q$ pigeons.  Let $\Phi$ be restriction from $S_M$ to
$S_L\times H$ followed by semisimplification in the $H$-variable: an $H$-module
decomposes into indecomposables, and semisimplification discards the summands of dimension
divisible by $p$ (the \emph{negligible} ones).  This is the OTI functor of \cite{hone2025}
for the semisimplification functor $\varphi_H$, which satisfies the hypotheses of their
Theorem~8.4.5 (stated there).  $\Phi$ is strong symmetric monoidal \cite[\S2.4]{hone2025}.  It
is \emph{not} exact on all $kS_M$-modules.

\paragraph{External input.}  Call $\lambda\vdash M$ \emph{$q$-stable} if
$\sum_{i\ge2}\lambda_i<q$ and, when $M<3q$, $\lambda_1\ge(M+q)/2$
\cite[Def.~6.1.1]{hone2025}.  Let $\mathrm{Pol}^{\mathrm{st}}_{M,q}$ be the Serre subcategory
of polynomial functors of degree $M$ whose simples are the $L_{\lambda^t}$ with $\lambda$
$q$-stable, and $\mathrm{Pol}^{\mathrm{tr}}_{L,q}$ its truncated counterpart in degree $L$, with
simples $L_{(\lambda-qe_1)^t}$.

\begin{theorem}[{\cite[Thm.~8.4.5]{hone2025}}]\label{thm:hone}
For $\varphi$ satisfying the two hypotheses there (semisimplification does), the OTI
functor corestricts to an equivalence of highest-weight categories
$\mathrm{Pol}^{\mathrm{st}}_{M,q}\to\mathrm{Pol}^{\mathrm{tr}}_{L,q}$.
\end{theorem}

By \cite[Lemma~7.3.2]{hone2025}, this functor commutes with Schur--Weyl duality: for the
Schur functor $\mathcal F$, $(\mathcal F\boxtimes1)\Phi\simeq\Phi\mathcal F$.  We use the
twisted Schur functor $\tilde{\mathcal F}=\mathcal F(-)\otimes\mathrm{sgn}$; twisting costs
nothing on the $H$-side, since $H$ is a $p$-group and its sign character is trivial in
characteristic $p$.

\paragraph{Step 1: exactness on the bounded-tail range.}  Let $\mathcal A_D(M)$ be the Serre
subcategory of polynomial functors whose simple labels are $\lambda^t$ with
$\lambda=(M-e,\tau)$, $|\tau|=e\le D$, and $\mathcal A_D(L)$ the same with $L$.  These
$\lambda$ are $q$-stable: $e\le D<q$, and $M-e\ge(M+q)/2$ is equivalent to $L\ge2e$.
Theorem~\ref{thm:hone} restricts to an exact equivalence
$\Theta\colon\mathcal A_D(M)\to\mathcal A_D(L)$ preserving $e$ and $\tau$.  The kernel
$K_D(M)$ of $\tilde{\mathcal F}$ on $\mathcal A_D(M)$ is generated by the simples whose
$\lambda$ is $p$-singular (some part repeated $p$ times).  Since
$M-e=L+q-e>e\ge\tau_1$, a repetition in $\lambda$ lies inside $\tau$ and remains one in
$(L-e,\tau)$, where $L-e\ge e\ge\tau_1$.  Hence $\Theta(K_D(M))\subseteq K_D(L)$.  (For
$L>2D$ this is an equality; at $L=2D$ it can be strict.)  The exact functor
$\tilde{\mathcal F}\Theta$ therefore factors through the Serre quotient
$\mathcal A_D(M)/K_D(M)$.

This quotient is the full image $\mathcal R_D(M)$ of $\tilde{\mathcal F}$ on
$\mathcal A_D(M)$, a full subcategory of $kS_M$-modules whose kernels, cokernels and
subobjects are the ordinary ones: for $X\in\mathcal A_D(M)$ and a submodule
$T\le\tilde{\mathcal F}X$, the pullback $X_T=X\times_{G\tilde{\mathcal F}X}GT$ along the
unit of the right adjoint $G$ of $\tilde{\mathcal F}$ is a subobject of $X$ with
$\tilde{\mathcal F}X_T=T$; lifting the graph of a module map in the same way identifies
morphisms.  (Compare \cite[Lemma~8.5.1]{hone2025}.)  By Lemma~7.3.2 the induced functor
$\mathcal R_D(M)\to\mathcal R_D(L)$ is $\Phi$.  So $\Phi$ is exact on $\mathcal R_D(M)$.  It
need not be faithful there, and we never use faithfulness.

\paragraph{Step 2: every object used lies in the range.}  The module $E_e(M)$ is
$\tilde{\mathcal F}$ of $\Lambda^{M-e}V\otimes V^{\otimes e}$, whose highest weight is
$(M-e,1^e)^t$; every composition factor has a label $\nu$ dominated by it, so $\nu^t$
dominates $(M-e,1^e)$, has first part $\ge M-e$ and tail of size $\le e$.  Hence $E_e(M)$,
$e\le D$, and so $F_d(M)$, $d\le D$, lie in $\mathcal R_D(M)$.  On a fixed hole set,
$U_j\otimes E_e\cong E_{e+1}\oplus E_e^{\oplus e}$ (the extra pigeon is new, or equals the
pigeon at one of the $e$ holes), so $U_j\otimes F_{D-1}$ lies in the range too, and with it
$U_j\otimes S$ for every submodule $S\le F_{D-1}$.  All intersections, images and sums in the
recursion therefore take place in $\mathcal R_D(M)$.

\paragraph{Step 3: the maps are the obvious ones.}  For a permutation module $k[X]$,
restricted to $H$, the $H$-orbits of size $>1$ give summands whose every $H$-endomorphism has
constant diagonal on the orbit and hence trace $0$; they are negligible and are killed.  So
$\Phi(k[X])=k[X^H]$, and an equivariant map between permutation modules goes to its block
from fixed basis vectors to fixed basis vectors.  A matching fixed by $H$ uses none of the
last $q$ pigeons (an edge at a moved pigeon would move), so $\Phi(F_d(M))=F_d(L)$ and
$\Phi(U_j(M))=U_j(L)$, compatibly with inclusions and with the monoidal structure.  On fixed
basis vectors: a pigeon axiom of a remaining pigeon is unchanged; a hole axiom
$\sum_{i}x_{ij}-1$ loses exactly the terms at moved pigeons, which form a non-trivial orbit,
and becomes the hole axiom of the smaller board; and $\mu_j$ sends $u_i\otimes x_N$ to the
same matching, the same new matching, or $0$, on both boards.

\paragraph{Step 4: the closure transfers.}  By induction on $t$, $\Phi$ carries
$C_t(M)\le F_D(M)$ to $C_t(L)\le F_D(L)$: exactness transfers the intersection with
$F_{D-1}$, the images and the sums; monoidality and Step~3 transfer
$\mu_j(U_j\otimes S_t)$.  Both chains stabilise after finitely many steps.

\paragraph{Step 5: detecting the constant.}  Let $c\colon k\to F_D(M)/C_\infty(M)$ be the
image of the constant.  If $c\ne0$ it is injective, since $k$ is simple; $\Phi$ is exact, so
$\Phi(c)$ is injective, and $\Phi(k)=k\ne0$, so $\Phi(c)\ne0$.  If $c=0$ then $\Phi(c)=0$.
By Step~4, $\Phi(c)$ is the corresponding map for $L$ pigeons.  So $1\in C_\infty(M)$ if and
only if $1\in C_\infty(L)$, and Lemma~\ref{lem:closure} finishes the step.  \qed

\subsection{The proof of part (a), through Harman's equivalences and tilting modules}\label{sec:periodicity-tilting}

For $m,m'>2D$, Theorem~\ref{thm:periodicity}(a) is proved without \cite{hone2025}; the proof is
given in full in Appendix~\ref{app:tilting}.  It transports the closure
along Harman's equivalence between the bounded-tail image categories of $S_m$ and $S_{m'}$ in weight
$\le D$~\cite[Thm.~2.9]{harman2015}.  The multiplication step is written as induction and
restriction of submodules, so the argument needs the equivalence to commute with induction
and restriction of arbitrary submodules, not only of permutation modules.  This is proved on
the Schur-algebra side, where induction and restriction become tensoring with the natural
module and a weight space of Levi restriction, both exact and preserving tilting modules; the
compatibility holds on tilting modules because the twisted Schur functor is fully faithful
there, and extends to all modules through $K^b(\mathrm{Tilt})\simeq D^b$.  The full
faithfulness is proved directly in every characteristic, because the identification Harman's
proof uses for it is false at $p=2$; the period must be a power of $p$ strictly greater than
$D$; and the strict inequality $m,m'>2D$ is needed (Section~\ref{sec:errata}).

\subsection{Checks}\label{sec:periodicity-checks}

The statements were tested by exact computation of $C_\infty$ (full row reduction over
$\mathbb F_p$, no use of symmetry).  For $D=2$ the predicate ``$1\in C_\infty$'' is periodic in
$m$ with period $4$ for $p=2$ ($m=5,\dots,12$), $3$ for $p=3$ ($m=5,\dots,13$) and $5$ for
$p=5$ ($m=5,\dots,14$), for $n=3,4$, and also for $n>m$.  Period $2$ fails for $p=2$, $D=2$
($n=3$: refutable at $m=5$, not at $m=7$), so the condition $q>D$ cannot be weakened to
$q\ge D$.  At the boundary $m=2D$ over $\mathbb F_2$ ($D=1,2,3$, $n\le7$) every congruent
triple agreed.  The mechanism itself was tested: for $q=p$ and $H$ generated by a $p$-cycle,
the Jordan type of the generator on $C_\infty(M)$ and on $F_D(M)/C_\infty(M)$ has no blocks of
size strictly between $1$ and $p$, and the number of blocks of size $1$ equals the dimension
of the corresponding space for $L=M-p$ pigeons (for example $(p,M,n,D)=(3,9,3,2)$:
$\dim C_\infty(9)=199$, with $88$ blocks of size $1$, and $\dim C_\infty(6)=88$).  These are
finite checks, not evidence for the general statement beyond the absence of a counterexample
at these sizes.

\paragraph{Who did what.}  All in run 20261104-0759; all reviews are model reviews.  The
closure formulation (Lemma~\ref{lem:closure}), the computations of
Section~\ref{sec:periodicity-checks} and the proof of Appendix~\ref{app:tilting} are by a Claude
agent (c1).  A Claude referee with no contact with the run reviewed that proof and found the
characteristic-$2$ error in \cite{harman2015}; a Codex agent (x2) reviewed it and supplied the
two-way scalar extension.  The proof of Section~\ref{sec:periodicity-oti} is by x2, who also found
the counterexample to \cite[Thm.~8.5.3]{hone2025} at the boundary and the exact-quotient argument
that avoids it; c1 reviewed it, and a second Codex agent (x1) adjudicated it, including the
boundary range.  After the run, a model referee re-read \cite{harman2015} and \cite{hone2025}
and re-ran the closure computations with its own code, and Astra adjudicated both proofs.  Astra
then reconstructed the proof of \cite[Thm.~8.4.5]{hone2025} and Lemma~7.3.2 from the source:
Lemma~7.3.2 holds, the proof of Theorem~8.4.5 has the gap described in
Section~\ref{sec:errata}, and Steps~1--5 hold conditional on that theorem.  A second model
re-checked the counterexample to Lemma~8.2.1 by hand.
